\documentclass[12pt,a4paper]{article}

% ── Bibliographie (doit précéder babel) ──────────────────────────
\usepackage{natbib}

% ── Encodage et langue ──────────────────────────────────────────
\usepackage[utf8]{inputenc}
\usepackage[T1]{fontenc}
\usepackage[french]{babel}

% ── Mise en page ─────────────────────────────────────────────────
\usepackage[top=2.5cm, bottom=2.5cm, left=2.5cm, right=2.5cm]{geometry}
\usepackage{setspace}
\onehalfspacing

% ── Mathématiques ────────────────────────────────────────────────
\usepackage{amsmath, amssymb}
\usepackage{mathtools}

% ── Tableaux ─────────────────────────────────────────────────────
\usepackage{booktabs}
\usepackage{array}
\usepackage{longtable}
\usepackage{multirow}
\usepackage{tabularx}

% ── Graphiques ───────────────────────────────────────────────────
\usepackage{graphicx}
\usepackage{caption}
\usepackage{subcaption}
\usepackage{float}

% ── Typographie ──────────────────────────────────────────────────
\usepackage{microtype}
\usepackage{lmodern}
\usepackage{parskip}
\setlength{\parindent}{0pt}
\setlength{\parskip}{6pt}

% ── Mise en page paysage ─────────────────────────────────────────
\usepackage{pdflscape}

% ── Liens ────────────────────────────────────────────────────────
\usepackage[colorlinks=true, linkcolor=black, citecolor=black,
            urlcolor=black]{hyperref}

% ── En-tête / pied de page ───────────────────────────────────────
\usepackage{fancyhdr}
\pagestyle{fancy}
\fancyhf{}
\fancyhead[L]{\small\textit{Système de visualisation  LOF et BiVAT}}
\fancyhead[R]{\small\textit{MOUKOKO EKAMBI G.M.}}
\fancyfoot[C]{\thepage}
\renewcommand{\headrulewidth}{0.4pt}

% ── Sections ─────────────────────────────────────────────────────
\usepackage{titlesec}
\titleformat{\section}{\large\bfseries}{\thesection.}{0.5em}{}
\titleformat{\subsection}{\normalsize\bfseries}{\thesubsection}{0.5em}{}
\titleformat{\subsubsection}{\normalsize\itshape}{\thesubsubsection}{0.5em}{}

% ── Commandes utilitaires ────────────────────────────────────────
\newcommand{\nd}{--}
\newcommand{\placeholder}[2]{%
  \begin{center}
  \fbox{\begin{minipage}{0.85\textwidth}
  \centering\vspace{1.2cm}
  \texttt{[PLACEHOLDER  #1]}\\[0.4em]
  \small\textit{#2}
  \vspace{1.2cm}
  \end{minipage}}
  \end{center}
}



% ════════════════════════════════════════════════════════════════
\begin{document}

% ── Page de titre ────────────────────────────────────────────────
\begin{titlepage}
\centering
\vspace*{2cm}

{\LARGE\bfseries
SYSTÈME DE VISUALISATION POUR LA DÉTECTION DE POINTS D'INTÉRÊT\\[0.4em]
\large\normalfont
Catalogue raisonné des graphiques, positionnement dans le workflow
et guide d'utilisation opérationnelle  Modèles LOF et BiVAT\\[0.4em]
}

\vspace{2cm}

{\large\bfseries MOUKOKO EKAMBI GEORGES MIGUEL}\\[0.5em]
{\normalsize\itshape
Stagiaire, Banque des États de l'Afrique Centrale (BEAC)\\
M2 Data Science \& IA  ENSPD / Université de Douala, LIDIA Laboratory
}

\vspace{2cm}

{\normalsize Août 2026}

\vfill

{\small
Document rédigé dans le cadre du stage de recherche à la Banque des États
de l'Afrique Centrale (BEAC), M2 Data Science \& IA,
ENSPD / Université de Douala  LIDIA Laboratory.
}
\end{titlepage}

% ── Résumé ───────────────────────────────────────────────────────
\begin{abstract}
Ce document décrit le système complet de visualisation associé à l'outil
de détection de points d'intérêt dans les données monétaires BEAC-CEMAC,
développé dans le cadre d'une étude comparative à deux modèles : le Local
Outlier Factor (LOF, Modèle~1) et BiVAT (\textit{Bidirectional
VAE-Associative Transformer}, Modèle~2). Pour chaque graphique, trois
informations sont documentées avec le même niveau de rigueur que les
documents méthodologiques compagnons : la pertinence analytique dans le
contexte BEAC, le positionnement précis dans le workflow de détection en
quatre phases (diagnostic des données, calibration interne, production des
résultats, comparaison inter-modèles), et les données requises pour sa
production. Le document identifie quatorze graphiques issus du corpus
initial, évalue leur pertinence individuelle, en écarte deux pour des
raisons méthodologiques documentées, et propose sept graphiques
supplémentaires absents du corpus initial mais nécessaires à la complétude
du système visuel. Le catalogue récapitulatif final (section~6) constitue
le guide d'implémentation opérationnelle de l'outil.

\bigskip
\noindent\textbf{Mots-clés :} visualisation, détection d'anomalies,
séries temporelles financières, LOF, BiVAT, BEAC, CEMAC, workflow,
points d'intérêt, audit macroprudentiel.
\end{abstract}

\tableofcontents
\newpage

% ════════════════════════════════════════════════════════════════
\section{Introduction et Architecture du Système Visuel}

\subsection{Pourquoi un système de visualisation dédié}

La détection de points d'intérêt dans les données monétaires BEAC-CEMAC
produit des scores numériques  scores LOF* et scores BiVAT  qui ne
sont exploitables opérationnellement que s'ils sont accompagnés d'un système
de visualisation structuré. Un superviseur prudentiel BEAC confronté à une
liste de scores numériques ne peut pas prendre de décision d'investigation
sans savoir à quoi correspondent ces scores dans le temps, quels indicateurs
IFS les portent, et si les deux modèles s'accordent sur les mêmes alertes.

Ce document traite le système de visualisation comme une composante
méthodologique à part entière, et non comme un simple habillage graphique
des résultats. Chaque graphique est analysé sous trois angles : sa pertinence
analytique dans le contexte BEAC (est-ce qu'il apporte une information
non disponible autrement ?), son positionnement dans le workflow de détection
(à quel moment du processus est-il produit et consulté ?), et ses dépendances
de calcul (quelles données et quels résultats intermédiaires sont nécessaires
à sa production ?).

\subsection{Architecture du workflow en quatre phases}

Le workflow de détection de points d'intérêt est organisé en quatre phases
séquentielles. Chaque phase a un objectif distinct et produit un ensemble
de graphiques spécifiques. La figure~\ref{fig:workflow} schématise cette
architecture.

\begin{figure}[H]
\centering
\placeholder{Fig. Workflow}{
Diagramme en quatre blocs séquentiels : Phase~1 Diagnostic des données
$\rightarrow$ Phase~2 Calibration et validation interne $\rightarrow$
Phase~3 Production et interprétation des résultats $\rightarrow$
Phase~4 Comparaison inter-modèles et décision finale. Chaque bloc liste
les graphiques associés et les données en entrée/sortie.
}
\caption{Architecture du workflow de détection en quatre phases.
Les graphiques associés à chaque phase sont indiqués entre parenthèses.}
\label{fig:workflow}
\end{figure}

\textbf{Phase~1  Diagnostic des données.} Objectif : valider la qualité
des données brutes et la conformité de chaque étape du prétraitement commun
avant d'engager les modèles. Cette phase est exécutée une seule fois, en
amont des deux modèles. Elle répond à la question : les données sont-elles
exploitables et le prétraitement produit-il les transformations attendues ?

\textbf{Phase~2  Calibration et validation interne.} Objectif : régler
les hyperparamètres de chaque modèle et valider que les deux modèles
fonctionnent correctement sur les données BEAC avant de produire des
résultats sur la fenêtre de test. Cette phase est exécutée séparément pour
LOF et pour BiVAT, sur la fenêtre d'entraînement et de calibration
(décembre~2001 -- décembre~2019).

\textbf{Phase~3  Production et interprétation des résultats.} Objectif :
appliquer les modèles calibrés sur la fenêtre de test (mars~2020 --
mars~2026) et produire les alertes avec leur interprétation économique.
C'est la phase principale de l'outil opérationnel. Les graphiques de cette
phase sont ceux qu'un superviseur BEAC consultera directement.

\textbf{Phase~4  Comparaison inter-modèles et décision finale.} Objectif :
confronter les résultats LOF et BiVAT, identifier les concordances et les
divergences, et formuler une recommandation opérationnelle structurée pour
la surveillance macroprudentielle.

\subsection{Convention de notation}

Dans ce document, les graphiques issus du corpus initial proposé sont
numérotés Fig.~1 à Fig.~14. Les graphiques supplémentaires identifiés
comme manquants sont notés Fig.~A à Fig.~G. La colonne \textbf{Priorité}
dans le catalogue récapitulatif (section~6) utilise trois niveaux :
\textbf{Essentiel} (nécessaire à la validité du workflow), \textbf{Important}
(apporte une information analytique significative) et \textbf{Complémentaire}
(utile pour l'exploration ou la présentation mais non critique).

% ════════════════════════════════════════════════════════════════
\section{Phase 1  Diagnostic des Données}

Cette phase regroupe les graphiques produits après le prétraitement commun
et avant l'exécution des modèles. Ils valident que les données sont
exploitables, que l'imputation est adéquate, que la décomposition STL
produit des résidus stationnaires, et que la structure de covariance des
résidus justifie la réduction de dimension.

\subsection{Fig.~1  Carte des données manquantes}

\subsubsection{Description}

Matrice binaire de dimensions (indicateurs IFS $\times$ mois), où chaque
case est noire si la valeur est manquante ou contient une erreur
\texttt{\#VALUE!} Excel, et blanche si la valeur est disponible. Une
colonne de droite indique le taux de complétude par indicateur ; une ligne
du bas indique le taux de complétude par mois.

\placeholder{Fig. 1  Carte des données manquantes}{
Matrice binaire (55 indicateurs × 195 mois). Axe vertical : indicateurs
IFS triés par taux de complétude décroissant. Axe horizontal : mois de
déc. 2001 à mars 2026. Cases noires = NaN ou \#VALUE!. Barres latérales
de taux de complétude. Fichier : \texttt{fig1\_missing\_data\_map.pdf}
}

\subsubsection{Pertinence analytique}

Ce graphique est le premier produit dans le workflow et l'un des plus
importants. Il remplit trois fonctions simultanées. D'une part, il rend
visible la répartition spatiale et temporelle des lacunes, permettant de
distinguer immédiatement les trois régimes d'imputation : lacunes courtes
isolées (interpolation linéaire), lacunes étendues internes (MICE), et
lacunes structurelles en début de série (troncature). D'autre part, il
révèle si des indicateurs entiers sont systématiquement manquants sur
certaines périodes  ce qui signalerait un problème de collecte
structurel à signaler à la BEAC indépendamment de toute anomalie de
contenu. Enfin, il documente empiriquement la justification de l'exclusion
des séries dont le taux d'imputation dépasse 20\,\%, évitant que cette
décision apparaisse arbitraire dans le rapport final.

\subsubsection{Positionnement dans le workflow}

\textbf{Phase~1  Diagnostic des données.} Ce graphique est produit
immédiatement après le chargement des données brutes, avant toute
imputation. Il constitue le point de départ de la validation de qualité.
Il est consulté une seule fois et ne sera pas mis à jour lors des exécutions
suivantes sur les mêmes fichiers sources.

\subsubsection{Données requises}

Données brutes au format XLSX, après lecture et conversion des erreurs
\texttt{\#VALUE!} en \texttt{NaN}. Aucune imputation ni normalisation
préalable. Un graphique par fichier pays/volet (12 graphiques au total)
ou une version agrégée multi-pays selon l'usage.

\subsubsection{Décision supportée}

Validation de la stratégie d'imputation (section~3.1 du document de
prétraitement). Identification des séries à exclure pour taux
d'imputation > 20\,\%. Détection de problèmes de collecte structurels
à signaler indépendamment de la détection d'anomalies.

\textbf{Priorité : Essentiel.}

\subsection{Fig.~2  Heatmap des corrélations des résidus STL}

\subsubsection{Description}

Matrice de corrélations de Pearson de dimensions ($55 \times 55$), calculée
sur les résidus $R_{i,t}$ après décomposition RobustSTL. Les indicateurs
sont réordonnés par clustering hiérarchique pour faire apparaître les blocs
de corrélation. L'échelle de couleur va du bleu foncé (corrélation $-1$)
au rouge foncé (corrélation $+1$), avec le blanc à zéro.

\placeholder{Fig. 2  Heatmap des corrélations des résidus STL}{
Matrice 55×55. Clustering hiérarchique Ward sur les lignes et colonnes.
Dendrogrammes latéraux. Diagonale = 1 (rouge foncé). Blocs visibles
attendus : sous-postes Crédits, sous-postes Dépôts, Titres.
Fichier : \texttt{fig2\_heatmap\_correlations\_stl.pdf}
}

\subsubsection{Pertinence analytique}

Ce graphique est produit sur les résidus STL, non sur les données brutes.
Cette distinction est fondamentale : sur données brutes, la heatmap
reflèterait principalement la tendance macroéconomique commune à tous les
indicateurs (les crédits et les dépôts croissent tous sur 2001--2026),
masquant les corrélations structurelles entre résidus. Sur résidus STL,
elle révèle les blocs de redondance entre sous-postes du bilan  par
exemple, les résidus des crédits aux ménages et aux sociétés non financières
corrèlent probablement fortement, justifiant leur compression en composantes
principales.

Ce graphique remplit deux fonctions complémentaires. Il justifie la
réduction de dimension par RPCA : si la heatmap révèle une structure
fortement bloc-diagonale, la variance est effectivement concentrée dans un
sous-espace de faible dimension, validant l'approche RPCA. Il justifie
également l'Association Discrepancy de BiVAT : les blocs de corrélation
visibles en Phase~1 deviennent les \textit{patterns d'association normaux}
que BiVAT apprend à modéliser, et dont les ruptures constituent des anomalies
collectives au sens de la lacune~L1.

\subsubsection{Positionnement dans le workflow}

\textbf{Phase~1  Diagnostic des données.} Ce graphique est produit
après la décomposition RobustSTL et avant la réduction de dimension. Il
est consulté une fois pour valider la structure des données d'entrée des
deux modèles.

\subsubsection{Données requises}

Résidus $R_{i,t}$ après décomposition RobustSTL et normalisation
médiane/MAD. Matrice de dimension $(N_{\mathrm{train}} \times 55)$ sur
la fenêtre d'entraînement uniquement (décembre~2001 -- novembre~2018).

\subsubsection{Décision supportée}

Justification de la réduction de dimension (nombre de blocs visibles
informe sur $k$). Identification des sous-espaces de corrélation pertinents
pour l'interprétation économique des anomalies. Validation que les résidus
STL ont une structure de covariance exploitable.

\textbf{Priorité : Essentiel.}

\subsection{Fig.~3  Décomposition RobustSTL à quatre panneaux}

\subsubsection{Description}

Pour trois indicateurs représentatifs choisis pour leur diversité de
comportement (un indicateur à forte tendance comme les Crédits totaux, un
à forte saisonnalité comme les Dépôts Transférables, un à forte volatilité
résiduelle), quatre panneaux empilés verticalement : série brute $x_{i,t}$,
tendance $T_{i,t}$, composante saisonnière $S_{i,t}$, résidu $R_{i,t}$.
L'axe des abscisses est commun (mois, 2001--2026). Les ruptures
macroéconomiques documentées (crise pétrolière 2015, COVID-19 2020) sont
annotées par des lignes verticales.

\placeholder{Fig. 3  Décomposition RobustSTL à quatre panneaux}{
3 sous-figures côte à côte (indicateur 1, 2, 3), chacune avec 4 panneaux
empilés. Hauteur totale : 24cm. Lignes verticales annotées : déc. 2015
(nadir crise pétrolière), mars 2020 (COVID-19).
Fichier : \texttt{fig3\_decomposition\_robststl.pdf}
}

\subsubsection{Pertinence analytique}

C'est la justification visuelle centrale du choix de travailler sur les
résidus plutôt que sur les données brutes, pour les deux modèles. Ce
graphique démontre de façon immédiatement accessible ce que le texte
méthodologique argumente : la tendance capte la croissance nominale sur
25 ans, la saisonnalité capte les cycles budgétaires annuels, et le résidu
 après extraction de ces deux composantes  est une série approximativement
stationnaire centrée sur zéro dont les déviations sont les candidats naturels
aux anomalies. Sans ce graphique, le lecteur doit accepter sur parole que
RobustSTL produit des résidus adaptés à la détection d'anomalies.

Le choix de trois indicateurs de comportements différents couvre le spectre
complet des données BEAC et permet au lecteur de généraliser visuellement
à l'ensemble des 55 séries.

\subsubsection{Positionnement dans le workflow}

\textbf{Phase~1  Diagnostic des données.} Ce graphique est produit
immédiatement après la décomposition RobustSTL, avant la normalisation et
la réduction de dimension. Il valide que la décomposition produit les
composantes attendues pour chaque type de série.

\subsubsection{Données requises}

Données brutes $x_{i,t}$ et composantes décomposées $T_{i,t}$, $S_{i,t}$,
$R_{i,t}$ après RobustSTL. Période complète 2001--2026. Sélection manuelle
des trois indicateurs représentatifs.

\subsubsection{Décision supportée}

Validation de la décomposition RobustSTL. Justification du travail sur
les résidus pour les deux modèles. Vérification visuelle que les anomalies
documentées (COVID-19 mars 2020) apparaissent comme des pics dans les
résidus et non dans la tendance.

\textbf{Priorité : Essentiel.}

\subsection{Fig.~A  Évolution de la variance expliquée par RPCA dans le
temps}

\subsubsection{Description}

Graphique en ligne montrant l'évolution du pourcentage de variance expliquée
par les $k = 5$ premières composantes RPCA, calculé sur des fenêtres
glissantes de 36 mois (pas de 1 mois) sur la période 2001--2026. L'axe
des ordonnées va de 0 à 100\,\%. Les ruptures documentées sont annotées.
Une ligne horizontale indique le seuil cible de 90\,\%.

\placeholder{Fig. A  Variance expliquée RPCA dans le temps}{
Courbe de la variance expliquée cumulée des 5 premières composantes RPCA
sur fenêtres glissantes de 36 mois. Annotations : crise pétrolière 2015,
COVID-19 2020. Ligne horizontale à 90\%. Si la courbe chute sous 90\% en
2020, c'est un signal de rupture de régime documenté.
Fichier : \texttt{figA\_variance\_rpca\_temporelle.pdf}
}

\subsubsection{Pertinence analytique}

Ce graphique est absent du corpus initial mais constitue une contribution
analytique majeure. Un scree plot global calculé sur l'ensemble de la
période 2001--2026 masque les changements de structure de covariance dans
le temps. Si la variance expliquée par les cinq premières composantes chute
significativement en 2020 (par exemple de 85\,\% à 65\,\%), cela signifie
que la structure de covariance entre indicateurs a changé  information
directement pertinente pour deux raisons. D'une part, elle signale une
rupture de régime indépendamment de tout score LOF ou BiVAT, constituant
une détection précoce de niveau structurel. D'autre part, elle documente
empiriquement l'une des spécificités des données financières BEAC (non-stationnarité
de la structure de covariance) qui justifie RevIN dans le pipeline BiVAT
et la Conformal Prediction pondérée.

\subsubsection{Positionnement dans le workflow}

\textbf{Phase~1  Diagnostic des données.} Ce graphique est produit après
la décomposition RobustSTL et la normalisation MAD, en parallèle de Fig.~2.
Il est consulté une fois pour établir si le nombre $k$ de composantes doit
être fixe ou adaptatif dans le temps.

\subsubsection{Données requises}

Résidus normalisés $x_{i,t}^*$ après RobustSTL et médiane/MAD. Application
de RPCA sur fenêtres glissantes de 36 mois. Période complète 2001--2026.

\subsubsection{Décision supportée}

Validation du choix de $k$ fixe vs adaptatif pour la RPCA du pipeline LOF.
Détection précoce de ruptures de régime au niveau structural (indépendante
des scores LOF et BiVAT). Justification empirique de RevIN et CP pondérée
pour BiVAT sur les périodes à forte chute de variance expliquée.

\textbf{Priorité : Important.}
% ════════════════════════════════════════════════════════════════
\section{Phase 2  Calibration et Validation Interne des Modèles}

Cette phase regroupe les graphiques produits pendant l'entraînement et la
calibration des deux modèles, sur la fenêtre décembre~2001 -- décembre~2019.
Ils valident que chaque modèle fonctionne correctement sur les données BEAC
avant d'être appliqué à la fenêtre de test. Aucun graphique de cette phase
n'est présenté dans le rapport final à la BEAC  ils constituent des
outils de validation interne pour le développeur et le chercheur.

\subsection{Fig.~4  Scree plot RPCA}

\subsubsection{Description}

Graphique à double axe : barres verticales représentant la variance
expliquée par chaque composante individuelle (axe gauche, pourcentage),
et courbe cumulée superposée (axe droit, pourcentage cumulé). Le nombre
de composantes retenues $k$ est indiqué par une ligne verticale pointillée
à l'intersection de la courbe cumulée avec le seuil de 90\,\%. L'axe des
abscisses va de la composante~1 à la composante~20.

\placeholder{Fig. 4  Scree plot RPCA}{
Barres : variance expliquée par composante (\%). Courbe : variance
cumulée (\%). Ligne horizontale à 90\%. Ligne verticale à $k$ retenu.
Annotation : ``$k = X$ composantes expliquent $Y$\% de la variance.''
Calculé sur résidus normalisés, fenêtre d'entraînement 2001--2018.
Fichier : \texttt{fig4\_scree\_plot\_rpca.pdf}
}

\subsubsection{Pertinence analytique}

Ce graphique justifie empiriquement le choix du nombre de composantes $k$
retenues pour la réduction de dimension du pipeline LOF. Sans lui, le
choix de $k$ est arbitraire. Le critère du coude (point d'inflexion de
la courbe cumulée) est visualisé directement. Si le coude apparaît à
$k = 8$ mais que 90\,\% de variance est atteint à $k = 12$, le graphique
permet de trancher entre les deux critères avec une justification visuelle
claire. Il révèle également le degré de redondance entre indicateurs : si
5 composantes expliquent 90\,\% de la variance pour 55 indicateurs, la
redondance est très élevée, ce qui est cohérent avec la structure des
sous-postes du bilan IFS.

Ce graphique est spécifique au pipeline LOF. BiVAT n'ayant pas de réduction
de dimension préalable, il n'a pas d'équivalent dans le pipeline BiVAT.

\subsubsection{Positionnement dans le workflow}

\textbf{Phase~2  Calibration LOF.} Ce graphique est produit lors de la
première application de RPCA sur la fenêtre d'entraînement. Il est consulté
une seule fois pour fixer $k$ avant d'exécuter le LOF. Il n'est pas
reproductible à chaque exécution opérationnelle  $k$ est fixé lors de
la calibration initiale.

\subsubsection{Données requises}

Résidus normalisés $x_{i,t}^*$ sur la fenêtre d'entraînement
(décembre~2001 -- novembre~2018). Application de RPCA. Valeurs propres
de la composante $L$ après décomposition.

\subsubsection{Décision supportée}

Détermination de $k$ pour la RPCA du pipeline LOF. Validation que la
réduction de dimension est justifiée (forte concentration de la variance).
Choix entre critère du coude et critère des 90\,\% si les deux divergent.

\textbf{Priorité : Essentiel (pipeline LOF uniquement).}

\subsection{Fig.~5  Biplot ACP (composantes 1 et 2)}

\subsubsection{Description}

Nuage de points des 195 observations projetées dans le plan des deux
premières composantes principales (après RPCA). Chaque point représente
un mois, coloré par année (dégradé chromatique de 2001 à 2026). Les
vecteurs des indicateurs dominants (ceux dont le loading dépasse 0.3 en
valeur absolue sur l'une des deux composantes) sont superposés. Les mois
identifiés comme candidats anomaux (scores LOF* élevés) sont marqués
d'un symbole distinct.

\placeholder{Fig. 5  Biplot ACP (composantes 1 et 2)}{
Nuage de 195 points colorés par année (gradient chaud 2001→2026 froid).
Vecteurs des indicateurs dominants en gris. Points anomaux candidats
marqués d'une croix rouge. Axes : CP1 ($X$\% variance), CP2 ($Y$\%
variance). Ellipse de confiance à 95\% sur les données d'entraînement.
Fichier : \texttt{fig5\_biplot\_acp.pdf}
}

\subsubsection{Pertinence analytique}

Ce graphique est un outil d'exploration, non de validation. Sa pertinence
est modérée car les deux premières composantes ne capturent généralement
qu'une fraction de la variance totale (typiquement 40 à 60\,\% selon la
structure des données BEAC), et les anomalies les plus intéressantes peuvent
vivre dans les composantes 3 à 5. Il est utile pour trois usages limités :
obtenir une première impression visuelle de la séparabilité temporelle des
données (les mois de 2020 sont-ils visiblement isolés des mois de 2010 ?) ;
identifier les indicateurs qui dominent les deux premières composantes ;
et vérifier qu'aucun artefact de calcul ne produit des projections aberrantes.

\textit{Réserve importante.} Ce graphique ne doit pas être utilisé comme
argument de validation LOF  les points visuellement isolés dans le plan
CP1/CP2 ne sont pas nécessairement ceux que le LOF identifie comme anomaux
dans l'espace complet à $k$ dimensions.

\subsubsection{Positionnement dans le workflow}

\textbf{Phase~2  Calibration LOF, usage exploratoire.} Ce graphique est
produit après le calcul de la RPCA et avant le calcul des scores LOF.
Il est consulté une seule fois à titre exploratoire. Il n'est pas inclus
dans le rapport opérationnel final à la BEAC.

\subsubsection{Données requises}

Composantes RPCA calculées sur la fenêtre d'entraînement. Projections des
195 observations sur CP1 et CP2. Loadings des 55 indicateurs.

\subsubsection{Décision supportée}

Exploration visuelle préliminaire. Identification des indicateurs dominants
dans l'espace réduit. Aucune décision méthodologique ne repose
exclusivement sur ce graphique.

\textbf{Priorité : Complémentaire.}

\subsection{Fig.~6  Projection UMAP}

\subsubsection{Description}

Projection bidimensionnelle des 195 observations par UMAP
(\textit{Uniform Manifold Approximation and Projection}), colorée par
score LOF* calculé sur la fenêtre complète. L'échelle de couleur va du
bleu (score LOF* faible, point normal) au rouge foncé (score LOF* élevé,
anomalie). Les mois des événements documentés (COVID-19, crise pétrolière)
sont annotés.

\placeholder{Fig. 6  Projection UMAP colorée par score LOF*}{
Nuage UMAP bidimensionnel. Colormap : bleu (LOF* bas) → rouge (LOF* élevé).
Points annotés : mars 2020, déc. 2015, et les 5 scores les plus élevés.
Paramètres UMAP : n\_neighbors=15, min\_dist=0.1, metric=euclidean.
Note : non-déterministe, fixer random\_state.
Fichier : \texttt{fig6\_projection\_umap.pdf}
}

\subsubsection{Pertinence analytique et réserves méthodologiques}

UMAP est un outil de visualisation exploratoire puissant mais son usage
dans ce contexte requiert deux réserves explicites documentées dans les
documents méthodologiques compagnons.

\textit{Réserve 1  Non-déterminisme.} UMAP produit des résultats
différents selon l'initialisation aléatoire. La graine aléatoire
(\texttt{random\_state}) doit être fixée et documentée pour garantir la
reproductibilité (contrainte C5 du pipeline LOF). Deux exécutions avec
des graines différentes produiront des projections visuellement différentes
même si les scores LOF* sont identiques.

\textit{Réserve 2  Distorsion des distances globales.} UMAP préserve
les structures locales mais distord les distances globales. Deux points
visuellement proches dans la projection UMAP peuvent être distants dans
l'espace à $k$ dimensions où le LOF calcule ses scores. Ce graphique
est une visualisation, non une validation. Il ne doit pas être présenté
comme une confirmation de la cohérence géométrique des scores LOF.

Dans ces limites, UMAP apporte une information visuelle utile que le biplot
ACP ne peut pas fournir : la capture des structures non-linéaires locales.
Si les mois de COVID-19 forment un cluster visuellement isolé dans la
projection UMAP, c'est une indication qualitative de leur atypicité 
mais seulement qualitative.

\subsubsection{Positionnement dans le workflow}

\textbf{Phase~2  Calibration LOF, usage exploratoire.} Ce graphique est
produit après le calcul des scores LOF* sur la fenêtre d'entraînement. Il
est consulté à titre exploratoire et placé en annexe du rapport interne,
avec la note méthodologique sur ses limites. Il n'est pas inclus dans le
rapport opérationnel final à la BEAC.

\subsubsection{Données requises}

Résidus normalisés $x_{i,t}^*$ sur la période complète (2001--2026).
Scores LOF* calculés. Paramètres UMAP fixés et documentés.

\subsubsection{Décision supportée}

Exploration visuelle non-linéaire des structures de données. Aucune décision
méthodologique ou opérationnelle ne repose sur ce graphique.

\textbf{Priorité : Complémentaire (annexe exploratoire uniquement).}

\subsection{Fig.~7  Profils LOF(MinPts)}

\subsubsection{Description}

Pour les 20 observations ayant reçu les scores LOF* les plus élevés, tracer
$\mathrm{LOF}_{\mathrm{MinPts}}(t)$ en fonction de MinPts pour
$\mathrm{MinPts} \in [10, 30]$. Chaque observation est une courbe distincte.
Une ligne horizontale indique le seuil IQR $\tau$. Les observations dont
la courbe reste au-dessus de $\tau$ sur toute la plage sont colorées en
rouge foncé (anomalies robustes) ; celles dont la courbe croise $\tau$
sont colorées en orange (anomalies sensibles).

\placeholder{Fig. 7  Profils LOF(MinPts) pour les top-20 anomalies}{
20 courbes, axe X : MinPts de 10 à 30, axe Y : score LOF(MinPts).
Courbes rouges : score > $\tau$ sur toute la plage (anomalies robustes).
Courbes oranges : score > $\tau$ seulement sur partie de la plage.
Ligne horizontale rouge pointillée : seuil $\tau$ (IQR rule).
Reproduction de la Fig. 8 de Breunig et al. (2000) sur données BEAC.
Fichier : \texttt{fig7\_profils\_lof\_minpts.pdf}
}

\subsubsection{Pertinence analytique}

C'est le graphique de validation méthodologique le plus important du
pipeline LOF. Il reproduit directement la Figure~8 de Breunig et al.\
(2000) sur les données BEAC et démontre empiriquement deux propriétés
fondamentales documentées dans le document méthodologique LOF. D'une part,
la non-monotonicité du LOF en fonction de MinPts : les courbes ne sont pas
monotones, confirmant que le score d'une même observation peut augmenter
ou diminuer selon le voisinage considéré. D'autre part, la robustesse des
anomalies les plus sévères : si les 5 à 10 anomalies les plus importantes
maintiennent un score élevé sur toute la plage $[10, 30]$, la stratégie
de la plage de valeurs est validée empiriquement sur ces données.

Ce graphique est également l'outil de calibration de la plage : si aucune
anomalie ne maintient un score élevé sur toute la plage, les bornes doivent
être révisées.

\subsubsection{Positionnement dans le workflow}

\textbf{Phase~2  Calibration LOF.} Ce graphique est produit lors de la
calibration de MinPts, sur la fenêtre d'entraînement. Il valide la plage
$[10, 30]$ retenue et identifie les anomalies robustes vs sensibles avant
l'application sur la fenêtre de test. Il est inclus dans le rapport interne
de validation, pas dans le rapport opérationnel BEAC.

\subsubsection{Données requises}

Calcul de $\mathrm{LOF}_{\mathrm{MinPts}}(t)$ pour chaque valeur entière
de MinPts dans $[10, 30]$ sur la fenêtre d'entraînement. Score maximum
$\mathrm{LOF}^*(t)$ et seuil $\tau$.

\subsubsection{Décision supportée}

Validation de la plage $[\mathrm{MinPts}_{\mathrm{LB}},
\mathrm{MinPts}_{\mathrm{UB}}]$. Distinction anomalies robustes (rouge) vs
anomalies sensibles (orange). Calibration du seuil $\tau$.

\textbf{Priorité : Essentiel (pipeline LOF).}

\subsection{Fig.~8  Distribution empirique des scores LOF*}

\subsubsection{Description}

Histogramme des scores $\mathrm{LOF}^*(t)$ sur la fenêtre d'entraînement,
avec une estimation par noyau (KDE) superposée. Trois lignes verticales
annotées : médiane, seuil IQR ($\tau_{\mathrm{IQR}}$) et seuil au
95\textsuperscript{e} percentile ($\tau_{95}$). Les scores dépassant
$\tau_{\mathrm{IQR}}$ sont colorés en orange ; ceux dépassant les deux
seuils simultanément en rouge (anomalies robustes).

\placeholder{Fig. 8  Distribution empirique des scores LOF*}{
Histogramme (bins=30) + KDE. Ligne verte : médiane. Ligne orange :
$\tau_{\mathrm{IQR}} = Q_3 + 1.5 \times \mathrm{IQR}$. Ligne rouge :
$\tau_{95}$ (percentile 95). Zone orange : $\tau_{\mathrm{IQR}} < \mathrm{LOF*}
< \tau_{95}$ (anomalies marginales). Zone rouge : $\mathrm{LOF*} > \tau_{95}$
(anomalies robustes). Annotation : nombre de points dans chaque zone.
Fichier : \texttt{fig8\_distribution\_lof.pdf}
}

\subsubsection{Pertinence analytique}

Ce graphique remplit deux fonctions. D'une part, il justifie visuellement
les deux seuils retenus en montrant leur position dans la distribution
empirique réelle des scores : si $\tau_{\mathrm{IQR}}$ et $\tau_{95}$
coïncident ou sont très proches, les deux critères convergent et les
anomalies robustes sont clairement délimitées. Si les deux seuils divergent
fortement, la zone intermédiaire est large, indiquant une population
d'anomalies marginales à examiner avec expertise domaine. D'autre part,
la forme de la distribution elle-même est informative : une distribution
unimodale indique une population homogène de comportements normaux avec
quelques outliers ; une distribution bimodale suggère deux sous-populations
distinctes (par exemple régime pré et post-COVID), information utile pour
la surveillance macroprudentielle.

\subsubsection{Positionnement dans le workflow}

\textbf{Phase~2  Calibration LOF.} Ce graphique est produit après le
calcul des scores LOF* sur la fenêtre d'entraînement, pour fixer les seuils
$\tau$ avant l'application sur la fenêtre de test. Les seuils calculés sur
l'entraînement sont appliqués sans recalcul sur le test.

\subsubsection{Données requises}

Scores $\mathrm{LOF}^*(t)$ sur la fenêtre d'entraînement
(décembre~2001 -- novembre~2018). Calcul de $Q_1$, $Q_3$, IQR, $\tau_{\mathrm{IQR}}$
et $\tau_{95}$.

\subsubsection{Décision supportée}

Fixation des seuils de détection $\tau_{\mathrm{IQR}}$ et $\tau_{95}$.
Évaluation de la proportion d'anomalies robustes vs marginales sur
l'entraînement. Détection d'éventuelles sous-populations dans les données
d'entraînement.

\textbf{Priorité : Essentiel (pipeline LOF).}

\subsection{Fig.~E  Intervalles de confiance BiVAT par Conformal
Prediction (calibration)}

\subsubsection{Description}

Graphique en ligne du score discriminant BiVAT
$s^{\mathrm{discrim}}_t$ sur la fenêtre de calibration
(décembre~2018 -- décembre~2019), avec les enveloppes de confiance CP à
90\,\% et 95\,\% superposées. Une ligne horizontale indique le quantile
empirique $\hat{q}_{1-\alpha}$ calculé sur cette fenêtre. Les mois dont
le score dépasse $\hat{q}_{1-\alpha}$ sont annotés.

\placeholder{Fig. E  Intervalles CP BiVAT (fenêtre de calibration)}{
Axe X : mois de déc. 2018 à déc. 2019 (13 mois). Axe Y : score
$s^{\mathrm{discrim}}_t$. Courbe centrale bleue : score BiVAT.
Enveloppe bleue claire : intervalle CP 90\%. Enveloppe bleue foncée :
intervalle CP 95\%. Ligne pointillée rouge : $\hat{q}_{0.9}$ (seuil
d'anomalie). Points au-dessus du seuil en rouge.
Fichier : \texttt{figE\_intervalles\_cp\_calibration.pdf}
}

\subsubsection{Pertinence analytique}

Ce graphique valide le bon fonctionnement de la Conformal Prediction sur
la fenêtre de calibration. Il permet de vérifier empiriquement que le
quantile $\hat{q}_{1-\alpha}$ est bien calibré : si exactement
$(1-\alpha) \times 100\,\%$ des observations de calibration sont en dessous
de $\hat{q}_{1-\alpha}$, la garantie formelle de coverage est validée sur
cet échantillon. Si le coverage observé s'écarte significativement du niveau
nominal, la Small Sample Beta Correction (SSBC) doit être ajustée.

Ce graphique répond également à la lacune~L3 du LOF de façon visuelle :
il montre ce qu'un score d'anomalie avec intervalles de confiance apporte
par rapport à un score scalaire déterministe  non seulement un rang
ordinal, mais une mesure de la certitude associée à chaque décision.

\subsubsection{Positionnement dans le workflow}

\textbf{Phase~2  Calibration BiVAT.} Ce graphique est produit lors de
la calibration de la Conformal Prediction sur la fenêtre de calibration.
Il valide le quantile $\hat{q}_{1-\alpha}$ avant application sur la fenêtre
de test. Inclus dans le rapport interne de validation.

\subsubsection{Données requises}

Scores BiVAT $s^{\mathrm{discrim}}_t$ sur la fenêtre de calibration.
Quantile empirique $\hat{q}_{1-\alpha}$ avec SSBC. Enveloppes CP à 90\,\%
et 95\,\%.

\subsubsection{Décision supportée}

Validation du coverage empirique de la CP sur la fenêtre de calibration.
Fixation du quantile $\hat{q}_{1-\alpha}$ pour la fenêtre de test.
Ajustement éventuel de la SSBC si le coverage observé diverge du niveau
nominal.

\textbf{Priorité : Essentiel (pipeline BiVAT).}

% ════════════════════════════════════════════════════════════════
\section{Phase 3  Production et Interprétation des Résultats}

Cette phase regroupe les graphiques produits sur la fenêtre de test
(mars~2020 -- mars~2026) après application des modèles calibrés. Ce sont
les graphiques qui constituent le rapport opérationnel final destiné aux
équipes de surveillance BEAC. Ils répondent à trois questions distinctes :
quand les anomalies se produisent-elles ? Quels indicateurs les portent ?
Sont-elles idiosyncratiques (un pays) ou systémiques (plusieurs pays) ?

\subsection{Fig.~9  Boxplot des scores LOF* par année}

\subsubsection{Description}

Boxplot des scores $\mathrm{LOF}^*(t)$ groupés par année civile, sur la
période complète 2001--2026. La ligne centrale de chaque boîte est la
médiane annuelle ; les bords de la boîte sont $Q_1$ et $Q_3$ ; les
moustaches s'étendent à $1{,}5 \times \mathrm{IQR}$ ; les points au-delà
sont des outliers de l'année considérée. Les lignes horizontales $\tau_{\mathrm{IQR}}$
et $\tau_{95}$ sont superposées pour référence. Les années contenant au
moins une anomalie robuste sont colorées en rouge.

\placeholder{Fig. 9  Boxplot des scores LOF* par année}{
25 boîtes (2001 à 2026). Axe X : année. Axe Y : score LOF*. Ligne rouge
pointillée : $\tau_{\mathrm{IQR}}$. Ligne orange pointillée : $\tau_{95}$.
Boîtes rouges : années avec au moins 1 anomalie robuste. Annotation :
nombre d'anomalies robustes par année au-dessus de chaque boîte rouge.
Fichier : \texttt{fig9\_boxplot\_lof\_annee.pdf}
}

\subsubsection{Pertinence analytique}

Ce graphique fournit une vue temporelle agrégée qui révèle si les anomalies
se concentrent dans des années spécifiques ou sont distribuées uniformément.
C'est la première validation externe qualitative : si les boîtes de 2020
et de 2015--2016 sont significativement plus hautes que les autres, la
détection LOF est cohérente avec les chocs macroéconomiques documentés.
Il révèle également la variabilité intra-annuelle : une boîte 2020 très
haute avec une grande dispersion indique que certains mois de 2020 sont
très anomaux et d'autres moins  information utile pour identifier le pic
du choc COVID-19 dans les données BEAC.

\subsubsection{Positionnement dans le workflow}

\textbf{Phase~3  Production des résultats LOF.} Ce graphique est produit
après le calcul des scores LOF* sur la fenêtre de test complète (mais peut
aussi incorporer la fenêtre d'entraînement pour une vue 2001--2026). Il
est inclus dans le rapport opérationnel comme vue d'ensemble temporelle.

\subsubsection{Données requises}

Scores $\mathrm{LOF}^*(t)$ sur la période complète 2001--2026.
Seuils $\tau_{\mathrm{IQR}}$ et $\tau_{95}$ calculés sur l'entraînement.

\subsubsection{Décision supportée}

Identification des années concentrant les anomalies. Première validation
externe qualitative par concordance avec les événements documentés.
Présentation synthétique pour un auditeur BEAC non spécialiste du ML.

\textbf{Priorité : Important.}

\subsection{Fig.~10  Série temporelle des scores LOF*(t)}

\subsubsection{Description}

Graphique en ligne du score $\mathrm{LOF}^*(t)$ sur l'axe temporel
complet 2001--2026. Deux lignes horizontales : $\tau_{\mathrm{IQR}}$
(orange) et $\tau_{95}$ (rouge). Les points au-dessus de $\tau_{95}$
(anomalies robustes) sont annotés avec leur date et l'événement
macroéconomique correspondant si disponible. La zone entre $\tau_{\mathrm{IQR}}$
et $\tau_{95}$ est colorée en orange pâle (anomalies marginales). La zone
au-dessus de $\tau_{95}$ est colorée en rouge pâle.

\placeholder{Fig. 10  Série temporelle LOF*(t)}{
Ligne continue bleue : score LOF*(t). Zone orange pâle : anomalies marginales.
Zone rouge pâle : anomalies robustes. Annotations : ``Mars 2020  COVID-19'',
``Déc. 2015  Nadir crise pétrolière'', et les 5 scores les plus élevés non
documentés avec point d'interrogation. Axe X : mois 2001--2026.
Axe Y : score LOF*.
Fichier : \texttt{fig10\_serie\_temporelle\_lof.pdf}
}

\subsubsection{Pertinence analytique}

C'est le graphique central du rapport opérationnel. Il synthétise
l'ensemble des résultats LOF en une visualisation unique lisible par un
superviseur BEAC sans formation en machine learning. Chaque pic correspond
à un mois potentiellement anomal ; les annotations font le lien avec les
événements économiques documentés. Les anomalies non annotées (pics sans
événement connu) sont les candidates les plus intéressantes : elles
représentent soit des événements non documentés dans la littérature
disponible, soit des erreurs de collecte, soit des artefacts du modèle 
trois hypothèses à investiguer par expertise domaine.

La présence de la zone intermédiaire (anomalies marginales entre les deux
seuils) est une innovation par rapport à un graphique seuil unique : elle
traduit visuellement la gradation de l'anomalie et la nécessité d'une
expertise différenciée selon le niveau de confiance.

\subsubsection{Positionnement dans le workflow}

\textbf{Phase~3  Production des résultats LOF.} C'est le premier
graphique consulté par un superviseur BEAC dans le rapport opérationnel.
Il est produit après le calcul des scores LOF* sur la fenêtre de test et
peut être mis à jour mensuellement lors d'une utilisation opérationnelle
continue.

\subsubsection{Données requises}

Scores $\mathrm{LOF}^*(t)$ sur la période complète 2001--2026.
Seuils $\tau_{\mathrm{IQR}}$ et $\tau_{95}$. Liste des événements
macroéconomiques documentés CEMAC avec leurs dates.

\subsubsection{Décision supportée}

Identification des mois anomaux et de leur sévérité. Première investigation
des pics non annotés. Mise à jour mensuelle de la surveillance. Rapport
opérationnel principal à la BEAC.

\textbf{Priorité : Essentiel.}

\subsection{Fig.~11  Superposition LOF + résidus STL + série brute}

\subsubsection{Description}

Pour chaque anomalie robuste identifiée, un graphique à trois panneaux
superposés sur la même fenêtre temporelle de $\pm 24$ mois autour du mois
anomal : panneau supérieur  série brute $x_{i^*,t}$ de l'indicateur
principal contributeur ; panneau médian  résidu STL $R_{i^*,t}$ du
même indicateur ; panneau inférieur  score $\mathrm{LOF}^*(t)$ avec
seuil $\tau$. Le mois anomal est indiqué par une ligne verticale rouge.

\placeholder{Fig. 11  Superposition LOF + résidu STL + série brute}{
3 panneaux pour chaque anomalie majeure. Fenêtre : $t-24$ à $t+24$ mois.
Panneau 1 (haut) : série brute $x_{i^*,t}$ en milliards FCFA.
Panneau 2 (milieu) : résidu $R_{i^*,t}$ (normalisé MAD).
Panneau 3 (bas) : score LOF*(t) avec lignes $\tau_{\mathrm{IQR}}$ et
$\tau_{95}$. Ligne rouge verticale : mois anomal.
Fichier : \texttt{fig11\_superposition\_lof\_indicateur\_[mois].pdf}
}

\subsubsection{Pertinence analytique}

Ce graphique démontre visuellement un point méthodologique central : le
LOF réagit à la déviation dans l'espace des résidus normalisés, pas au
niveau absolu de la série brute. En montrant simultanément les trois niveaux
de représentation (série brute, résidu STL, score LOF), il rend
pédagogiquement visible pourquoi le prétraitement STL est indispensable.
Un pic dans le résidu sans pic apparent dans la série brute correspond
typiquement à une anomalie contextuelle (le niveau absolu est normal mais
la dynamique est atypique)  précisément le type d'anomalie que le LOF
peut capturer mais que la simple règle des $n$ sigmas sur données brutes
raterait.

\textit{Correction par rapport au corpus initial.} Le corpus initial
proposait de superposer LOF avec la série brute uniquement. Superposer
avec les résidus STL est plus informatif et méthodologiquement cohérent :
c'est sur les résidus que le LOF opère.

\subsubsection{Positionnement dans le workflow}

\textbf{Phase~3  Interprétation des résultats LOF.} Ce graphique est
produit après identification des anomalies robustes, une instance par
anomalie. Il est inclus dans le rapport opérationnel comme annexe
d'interprétation pour chaque alerte.

\subsubsection{Données requises}

Scores LOF*(t), résidus STL $R_{i,t}$, séries brutes $x_{i,t}$.
Identification de l'indicateur principal contributeur (via Fig.~13 ou
Fig.~F de l'attribution SHAP BiVAT). Seuils $\tau$.

\subsubsection{Décision supportée}

Interprétation économique de chaque anomalie robuste. Distinction entre
anomalie de niveau (pic dans la série brute) et anomalie de dynamique
(pic dans le résidu sans pic apparent dans la brute). Documentation des
alertes pour l'audit COBAC.

\textbf{Priorité : Important.}

\subsection{Fig.~12  Bar chart horizontal des indicateurs contributeurs}

\subsubsection{Description}

Pour chaque mois anomal, un bar chart horizontal trié par contribution
décroissante, montrant les 15 indicateurs dont les résidus normalisés
sont les plus élevés en valeur absolue au mois $t$. Les barres positives
(déviation à la hausse) sont en rouge ; les barres négatives (déviation
à la baisse) sont en bleu. Les valeurs numériques sont affichées en bout
de barre.

\placeholder{Fig. 12  Bar chart des indicateurs contributeurs}{
Bar chart horizontal. 15 indicateurs, triés par $|R_{i,t}^*|$ décroissant.
Barres rouges : $R_{i,t}^* > 0$. Barres bleues : $R_{i,t}^* < 0$.
Valeurs numériques en bout de barre. Titre : ``Indicateurs contributeurs
 [Mois anomal]''. Un graphique par mois anomal.
Fichier : \texttt{fig12\_contributeurs\_[mois].pdf}
}

\subsubsection{Pertinence analytique et correction du corpus initial}

Le corpus initial proposait un radar chart pour cette fonction.
Le bar chart horizontal lui est préféré pour trois raisons documentées.
D'une part, la lisibilité : avec 10 à 15 indicateurs aux noms longs de
type \texttt{Crédits\_ENF\_MN}, un radar chart est illisible pour un
auditeur BEAC. D'autre part, la précision : un bar chart permet de lire
les valeurs numériques exactes, critère d'auditabilité requis par les
exigences COBAC. D'autre part enfin, la distinction positive/négative :
un radar chart ne distingue pas visuellement les déviations à la hausse
des déviations à la baisse, information économiquement fondamentale (une
contraction des crédits et une expansion des dépôts sont deux signaux
opposés).

Le radar chart peut être conservé en complément comme visualisation
synthétique pour les présentations, mais le bar chart est l'outil d'audit.

\subsubsection{Positionnement dans le workflow}

\textbf{Phase~3  Interprétation économique des résultats LOF.}
Ce graphique est produit pour chaque anomalie robuste, en complément de
Fig.~11. Il répond à la question : quels indicateurs du bilan portent
l'anomalie ? Il est inclus dans le rapport opérationnel comme outil
d'interprétation économique pour chaque alerte.

\subsubsection{Données requises}

Résidus normalisés $R_{i,t}^*$ au mois anomal $t$ pour les 55 indicateurs.
Tri par valeur absolue décroissante. Identification des 15 indicateurs
les plus actifs.

\subsubsection{Décision supportée}

Attribution de l'anomalie aux indicateurs IFS contributeurs. Distinction
entre déviations positives et négatives. Documentation pour l'audit COBAC.
Formulation de l'hypothèse économique (erreur de collecte vs événement
réel).

\textbf{Priorité : Essentiel.}

\subsection{Fig.~13  Heatmap des anomalies (indicateurs × mois anomaux)}

\subsubsection{Description}

Matrice de chaleur de dimensions ($55 \times N_{\mathrm{anomalies}}$), où
chaque colonne est un mois anomal identifié (score LOF* > $\tau_{\mathrm{IQR}}$)
et chaque ligne est un indicateur IFS. La couleur encode l'intensité du
résidu normalisé $R_{i,t}^*$ : rouge foncé (déviation positive forte),
bleu foncé (déviation négative forte), blanc (résidu nul). Les indicateurs
sont triés par clustering hiérarchique sur les colonnes anomales.

\placeholder{Fig. 13  Heatmap des anomalies (indicateurs × mois anomaux)}{
Matrice (55 lignes × $N$ colonnes). Colormap divergente :
bleu foncé (résidu $< -3$) → blanc (résidu $= 0$) → rouge foncé (résidu $> 3$).
Dendrogramme sur les lignes (clustering indicateurs). Annotations sur les
colonnes : date et événement économique. Barre latérale droite : famille
IFS (Crédits, Dépôts, Titres\ldots)
Fichier : \texttt{fig13\_heatmap\_anomalies.pdf}
}

\subsubsection{Pertinence analytique}

C'est le graphique d'interprétation économique le plus riche du système
visuel. Il répond simultanément à deux questions : quels mois sont anomaux
(colonnes) ET quels indicateurs portent la déviation à chaque mois anomal
(lignes). La lecture croisée est immédiate pour un auditeur BEAC : si les
colonnes de 2020 montrent des déviations rouges concentrées sur les lignes
\texttt{Crédits} et des déviations bleues sur les lignes \texttt{Dépôts},
le choc COVID-19 a simultanément contracté les crédits et gonflé les dépôts
de précaution  interprétation économique directement exploitable.

Le clustering hiérarchique sur les lignes fait apparaître naturellement les
blocs d'indicateurs co-anomaux, révélant la structure des chocs dans le bilan
2SR. Ce graphique satisfait directement les exigences d'auditabilité COBAC
documentées dans le document BiVAT (BIS FSI 2024, ECB 2025).

\subsubsection{Positionnement dans le workflow}

\textbf{Phase~3  Interprétation économique globale.} Ce graphique est
produit après identification de l'ensemble des mois anomaux sur la fenêtre
de test. Il est inclus dans le rapport opérationnel comme vue d'ensemble
de l'interprétation économique des anomalies.

\subsubsection{Données requises}

Résidus normalisés $R_{i,t}^*$ pour tous les mois anomaux identifiés et
tous les 55 indicateurs. Liste des mois anomaux (score LOF* > $\tau_{\mathrm{IQR}}$).

\subsubsection{Décision supportée}

Interprétation économique des anomalies par indicateur. Identification
des blocs d'indicateurs co-anomaux (clusters de lignes). Documentation
pour l'audit COBAC. Formulation des hypothèses économiques sur la nature
des chocs.

\textbf{Priorité : Essentiel.}

\subsection{Fig.~14  Carte de chaleur inter-pays}

\subsubsection{Description}

Matrice de chaleur de dimensions (mois × pays CEMAC), où chaque cellule
encode le score LOF* normalisé du pays $p$ au mois $t$, ramené sur une
échelle commune $[0, 1]$ pour permettre la comparaison inter-pays. Les
six pays (Cameroun, Congo, Gabon, Guinée équatoriale, RCA, Tchad) sont
en colonnes ; les mois 2001--2026 en lignes. Les cellules dépassant le
seuil normalisé correspondant à $\tau_{\mathrm{IQR}}$ sont encadrées
en noir. Un panneau auxiliaire à droite indique, pour chaque mois, le
nombre de pays dépassant simultanément leur seuil national.

\placeholder{Fig. 14  Carte de chaleur inter-pays (mois × pays CEMAC)}{
Matrice (195 lignes × 6 colonnes). Colormap : blanc (score normalisé 0) →
rouge foncé (score normalisé 1). Encadrés noirs : cellules > seuil normalisé.
Panneau droit : nombre de pays en alerte simultanée (barres grises).
Ligne pointillée rouge : seuil systémique ($m \geq 3$ pays simultanément).
Fichier : \texttt{fig14\_heatmap\_interpays.pdf}
}

\subsubsection{Pertinence analytique}

Ce graphique implémente visuellement le niveau~4 du protocole de validation
LOF documenté dans le document méthodologique LOF : identification des
anomalies systémiques CEMAC ($m \geq 3$ pays simultanément au-dessus de
leur seuil). C'est la distinction fondamentale entre une anomalie
idiosyncratique (un seul pays, vraisemblablement erreur de collecte ou
choc local) et une anomalie systémique (plusieurs pays, vraisemblablement
choc macroéconomique CEMAC). Le panneau auxiliaire droite rend cette
distinction immédiatement lisible sans calcul supplémentaire.

\subsubsection{Positionnement dans le workflow}

\textbf{Phase~3  Production des résultats LOF, validation inter-pays.}
Ce graphique est produit après le calcul des scores LOF* sur les 12 fichiers
(6 pays × 2 volets Actif/Passif). Il est inclus dans le rapport opérationnel
comme outil de détection des anomalies systémiques CEMAC.

\subsubsection{Données requises}

Scores LOF* pour les 6 pays (Actif ou Passif selon le volet analysé).
Normalisation inter-pays des scores sur $[0, 1]$ (min-max sur chaque pays
séparément). Seuils nationaux $\tau_p$ pour chaque pays.

\subsubsection{Décision supportée}

Identification des anomalies systémiques CEMAC ($m \geq 3$ pays). Distinction
anomalie idiosyncratique vs systémique. Hiérarchisation des alertes par
niveau de risque macroprudentiel. Transmission à la surveillance régionale
BEAC pour les anomalies systémiques.

\textbf{Priorité : Essentiel.}

\subsection{Fig.~B  Concordance Actif--Passif par pays}

\subsubsection{Description}

Pour chaque pays, un graphique à deux panneaux superposés : score LOF*
Actif (panneau supérieur) et score LOF* Passif (panneau inférieur), sur
le même axe temporel 2001--2026. Les deux panneaux partagent l'axe des
abscisses et le même seuil $\tau$. Les mois où les deux panneaux dépassent
simultanément $\tau$ sont colorés en rouge (anomalie concordante) ; les
mois où un seul panneau dépasse $\tau$ sont colorés en orange (anomalie
discordante).

\placeholder{Fig. B  Concordance Actif--Passif par pays}{
2 panneaux empilés. Panneau haut : LOF* Actif. Panneau bas : LOF* Passif.
Lignes $\tau_{\mathrm{IQR}}$ sur les deux panneaux. Zones rouges :
concordance (les deux > $\tau$). Zones oranges : discordance (un seul > $\tau$).
Annotation des concordances : ``Événement économique réel probable.''
Annotation des discordances : ``Vérification comptable requise.''
6 graphiques (un par pays CEMAC).
Fichier : \texttt{figB\_concordance\_actif\_passif\_[pays].pdf}
}

\subsubsection{Pertinence analytique}

Ce graphique implémente le niveau~5 du protocole de validation LOF
(cohérence Actif--Passif). L'identité comptable
$\mathrm{Total\;Actif} = \mathrm{Total\;Passif}$ constitue une contrainte
de validation interne au dataset BEAC. Une anomalie concordante (signal
sur les deux volets) indique un événement économique réel affectant la
structure du bilan dans son ensemble. Une anomalie discordante suggère
prioritairement une erreur de collecte ou un problème de réconciliation
comptable. Cette distinction est documentée dans le document de prétraitement
commun (section~2.1) et opérationnalisée ici visuellement.

Ce graphique est absent du corpus initial et constitue une contribution
originale du système visuel, directement liée à la structure multi-fichiers
du dataset BEAC (12 fichiers = 6 pays × 2 volets).

\subsubsection{Positionnement dans le workflow}

\textbf{Phase~3  Validation des résultats LOF, niveau Actif--Passif.}
Ce graphique est produit après le calcul des scores LOF* sur les 12 fichiers.
Il est inclus dans le rapport opérationnel comme outil de triage des alertes
(événement économique vs erreur de collecte).

\subsubsection{Données requises}

Scores LOF* Actif et LOF* Passif pour chaque pays, sur la période complète.
Seuils $\tau_{\mathrm{IQR}}$ et $\tau_{95}$ calculés séparément pour chaque
fichier.

\subsubsection{Décision supportée}

Triage des alertes entre événements économiques réels (anomalies concordantes)
et problèmes de collecte (anomalies discordantes). Hiérarchisation des
investigations : les anomalies concordantes ont priorité sur les discordantes.
Signalement des anomalies discordantes à l'équipe de collecte BEAC.

\textbf{Priorité : Essentiel.}
% ════════════════════════════════════════════════════════════════
\section{Phase 4  Comparaison Inter-modèles et Décision Finale}

Cette phase regroupe les graphiques produits après l'exécution des deux
modèles sur la fenêtre de test. Ils confrontent les résultats LOF et BiVAT,
identifient les concordances et les divergences, et produisent les éléments
visuels nécessaires à la décision opérationnelle finale. Ces graphiques
constituent la valeur ajoutée principale du système comparatif à deux
modèles par rapport à un déploiement mono-modèle.

\subsection{Fig.~A bis  Comparaison LOF vs BiVAT dans le temps}
\label{subsec:figD}

\subsubsection{Description}

Graphique à deux courbes superposées sur le même axe temporel
(mars~2020 -- mars~2026) : score $\mathrm{LOF}^*(t)$ normalisé sur $[0,1]$
(courbe bleue) et score discriminant BiVAT $s^{\mathrm{discrim}}_t$ normalisé
sur $[0,1]$ (courbe rouge). Les seuils normalisés correspondant à
$\tau_{\mathrm{IQR}}$ (LOF) et $\hat{q}_{1-\alpha}$ (BiVAT) sont tracés
en lignes pointillées de la couleur de chaque modèle. Quatre zones de
couleur sont définies : vert pâle (les deux modèles sous leur seuil 
accord sur la normalité), rouge pâle (les deux modèles au-dessus de leur
seuil  accord sur l'anomalie), orange gauche (LOF seul au-dessus 
anomalie ponctuelle non temporelle), orange droit (BiVAT seul au-dessus 
anomalie contextuelle ou collective non détectée par LOF).

\placeholder{Fig. D  Comparaison LOF vs BiVAT dans le temps}{
Deux courbes normalisées sur [0,1]. Courbe bleue : LOF*(t)/max(LOF*).
Courbe rouge : $s^{\mathrm{discrim}}_t$/max($s^{\mathrm{discrim}}$).
Zone verte pâle : accord normalité. Zone rouge pâle : accord anomalie.
Zone orange gauche : LOF seul. Zone orange droit : BiVAT seul.
Annotations : événements documentés. Axe X : mars 2020 -- mars 2026.
Fichier : \texttt{figD\_comparaison\_lof\_bivat.pdf}
}

\subsubsection{Pertinence analytique}

C'est le graphique de synthèse comparative centrale de l'étude. Il
implémente visuellement les métriques M1 (Precision@k), M2 (rang des
anomalies connues) et M3 (cohérence inter-modèles) documentées dans les
deux documents méthodologiques compagnons. Les quatre zones colorées
structurent la lecture de la comparaison.

La zone verte (accord sur la normalité) valide la robustesse des deux
modèles sur les périodes stables. La zone rouge (accord sur l'anomalie)
identifie les alertes de niveau~1  les plus fiables car confirmées par
deux approches radicalement différentes (densité locale vs reconstruction
probabiliste). La zone orange gauche (LOF seul) identifie les anomalies
ponctuelles détectées par LOF mais pas par BiVAT : ce sont typiquement des
déviations isolées dans l'espace des résidus sans structure temporelle 
potentiellement des erreurs de saisie. La zone orange droit (BiVAT seul)
identifie les anomalies contextuelles ou collectives que le LOF rate
structurellement en raison de la lacune~L1 : ce sont précisément les
anomalies les plus pertinentes pour la surveillance macroprudentielle.

Ce graphique rend visible empiriquement, sur les données BEAC réelles,
l'écart de $-33{,}98$ points de F1 entre LOF et les méthodes temporelles
documenté dans le document LOF (section~5.1, Table~\ref{tab:l1-benchmark}).

\subsubsection{Positionnement dans le workflow}

\textbf{Phase~4  Comparaison inter-modèles.} Ce graphique est produit
après l'application des deux modèles calibrés sur la fenêtre de test.
C'est le premier graphique de la Phase~4, point d'entrée de la comparaison.
Il est inclus dans le rapport opérationnel comme outil de décision finale.

\subsubsection{Données requises}

Scores $\mathrm{LOF}^*(t)$ sur la fenêtre de test. Scores BiVAT
$s^{\mathrm{discrim}}_t$ sur la fenêtre de test. Normalisation min-max
de chaque série de scores séparément. Seuils $\tau_{\mathrm{IQR}}$ (LOF)
et $\hat{q}_{1-\alpha}$ (BiVAT) normalisés.

\subsubsection{Décision supportée}

Classification des alertes en quatre niveaux : accord normalité, accord
anomalie (niveau~1), LOF seul (niveau~2a), BiVAT seul (niveau~2b).
Hiérarchisation des investigations. Validation externe par concordance
avec les événements documentés. Quantification de la métrique M3
(cohérence inter-modèles).

\textbf{Priorité : Essentiel.}

\subsection{Fig.~E bis  Intervalles de confiance BiVAT sur la fenêtre
de test}

\subsubsection{Description}

Graphique en ligne du score discriminant BiVAT $s^{\mathrm{discrim}}_t$
sur la fenêtre de test (mars~2020 -- mars~2026), avec les enveloppes de
confiance Conformal Prediction à 90\,\% et 95\,\% superposées. Le quantile
empirique $\hat{q}_{1-\alpha}$ (calculé sur le calibration set) est tracé
en ligne pointillée rouge. Les mois dont le score dépasse $\hat{q}_{1-\alpha}$
sont annotés avec leur date. Les intervalles CP sont calculés par rolling-origin
(section~6.3 du document de prétraitement commun). Une bande grisée indique
les périodes de rupture de régime documentées où la CP pondérée est activée.

\placeholder{Fig. E bis  Score BiVAT + intervalles CP (fenêtre de test)}{
Axe X : mars 2020 -- mars 2026. Axe Y : score $s^{\mathrm{discrim}}_t$.
Courbe rouge : score BiVAT. Enveloppe rouge claire : intervalle CP 90\%.
Enveloppe rouge foncé : intervalle CP 95\%. Ligne pointillée rouge :
$\hat{q}_{0.9}$. Points annotés : mois dépassant le seuil. Bande grise :
mars 2020 -- déc. 2021 (rupture COVID, CP pondérée active).
Fichier : \texttt{figE\_bis\_bivat\_cp\_test.pdf}
}

\subsubsection{Pertinence analytique}

Ce graphique constitue la réponse visuelle directe à la lacune~L3 du LOF
(absence de quantification d'incertitude). Il montre ce qu'un score
d'anomalie avec intervalles de confiance apporte par rapport au score scalaire
déterministe LOF* : pour chaque mois, non seulement un rang ordinal, mais
une mesure de la certitude associée à la décision. Un mois dont le score
BiVAT dépasse $\hat{q}_{1-\alpha}$ mais dont l'intervalle inférieur de CP
reste proche du seuil est une anomalie à surveillance  une anomalie dont
le score central est élevé mais dont l'incertitude est grande. Un mois dont
le score est très au-dessus du seuil ET dont l'intervalle inférieur dépasse
aussi le seuil est une anomalie robuste à signalement immédiat.

La bande grisée de CP pondérée sur la période COVID-19 rend visible
l'adaptation du système aux ruptures de régime : les intervalles s'élargissent
pendant cette période, reflétant fidèlement l'incertitude accrue du modèle.

\subsubsection{Positionnement dans le workflow}

\textbf{Phase~4  Production des résultats BiVAT.} Ce graphique est
produit après l'application de BiVAT calibré sur la fenêtre de test et
le calcul des intervalles CP par rolling-origin. Il est inclus dans le
rapport opérationnel comme graphique principal BiVAT, en parallèle de
Fig.~10 pour LOF.

\subsubsection{Données requises}

Scores BiVAT $s^{\mathrm{discrim}}_t$ sur la fenêtre de test. Quantile
CP $\hat{q}_{1-\alpha}$ calculé sur le calibration set avec SSBC. Enveloppes
rolling-origin CP à 90\,\% et 95\,\%. Indicateur de période de CP pondérée.

\subsubsection{Décision supportée}

Identification des mois anomaux BiVAT avec niveau de confiance associé.
Classification : anomalie certaine (score + intervalle inférieur > seuil)
vs anomalie ambiguë (score > seuil mais intervalle inférieur < seuil).
Rapport opérationnel BiVAT principal à la BEAC.

\textbf{Priorité : Essentiel.}

\subsection{Fig.~F  Attribution SHAP par indicateur pour les top-5
anomalies BiVAT}

\subsubsection{Description}

Pour les cinq mois les plus anomaux selon BiVAT (score
$s^{\mathrm{discrim}}_t$ le plus élevé sur la fenêtre de test), un bar
chart horizontal des valeurs DeepSHAP par indicateur IFS, trié par
contribution absolue décroissante, affichant les 15 indicateurs les plus
contributeurs. Les barres positives (contribution à la hausse du score
d'anomalie) sont en rouge ; les barres négatives (contribution à la baisse)
en bleu. Les valeurs SHAP normalisées sont exprimées en pourcentage de
contribution totale. Un sous-titre indique le mois, le score BiVAT et le
score LOF correspondant (pour comparaison inter-modèles).

\placeholder{Fig. F  Attribution SHAP par indicateur (top-5 BiVAT)}{
5 bar charts horizontaux (un par anomalie majeure BiVAT).
15 indicateurs par graphique, triés par $|\phi_d|$ décroissant.
Barres rouges : $\phi_d > 0$ (contribution positive à l'anomalie).
Barres bleues : $\phi_d < 0$ (contribution négative).
Valeurs en \% de contribution totale. Sous-titre : ``[Mois] 
Score BiVAT : X.XX  Score LOF : Y.YY''.
Fichier : \texttt{figF\_shap\_[mois].pdf}
}

\subsubsection{Pertinence analytique}

Ce graphique constitue la réponse visuelle directe à la lacune~L4 du LOF
(absence d'explicabilité par indicateur) et représente l'output opérationnel
final le plus distinctif de BiVAT par rapport au LOF. Il produit exactement
ce qu'exigent les standards réglementaires documentés dans le document BiVAT :
BIS FSI Paper (2024), ECB Guide to Internal Models (2025) et SHARC
(arXiv:2607.05484). Pour chaque anomalie majeure, un superviseur BEAC peut
lire directement : ``l'anomalie de mars~2020 est attribuée à 67\,\% à la
série \texttt{Crédits\_ENF} et à 22\,\% à \texttt{Dépôts\_transf}'' 
information non disponible depuis le LOF.

Le sous-titre indiquant le score LOF correspondant permet également de
comparer directement les deux modèles pour chaque anomalie : si BiVAT
attribue un score élevé à un mois et que LOF lui attribue un score faible,
le graphique SHAP montre pourquoi  BiVAT a détecté une anomalie
contextuelle dans la combinaison de plusieurs indicateurs, que le LOF
dans l'espace réduit RPCA n'a pas vue.

\subsubsection{Positionnement dans le workflow}

\textbf{Phase~4  Interprétation des résultats BiVAT et comparaison.}
Ce graphique est produit après identification des cinq anomalies les plus
sévères de BiVAT sur la fenêtre de test. Il est inclus dans le rapport
opérationnel comme annexe d'interprétation réglementaire pour chaque
alerte BiVAT majeure.

\subsubsection{Données requises}

Valeurs DeepSHAP $\phi_d(\mathbf{x}_t)$ calculées pour chaque indicateur
$d$ aux cinq mois anomaux les plus sévères. Scores BiVAT et LOF
correspondants. Normalisation des valeurs SHAP en pourcentage de contribution
totale.

\subsubsection{Décision supportée}

Attribution causale de chaque anomalie BiVAT aux indicateurs IFS
contributeurs. Documentation réglementaire conforme aux exigences
COBAC/ICAAP. Comparaison inter-modèles pour les anomalies BiVAT non
détectées par LOF. Formulation de l'hypothèse économique par indicateur.

\textbf{Priorité : Essentiel (pipeline BiVAT).}

\subsection{Fig.~G  Carte de confusion LOF/BiVAT par pays et type
d'anomalie}

\subsubsection{Description}

Tableau de chaleur de dimensions ($6 \times 4$) : six lignes pour les
pays CEMAC et quatre colonnes pour les types d'anomalie (ponctuelle,
contextuelle, collective, systémique). Chaque cellule encode le taux de
concordance LOF--BiVAT pour ce type d'anomalie dans ce pays, sur la fenêtre
de test. L'échelle va du blanc (concordance nulle ou type non détecté) au
vert foncé (concordance maximale). Les cellules où les deux modèles
s'accordent sur plus de 70\,\% des alertes sont encadrées.

La classification des anomalies par type repose sur : ponctuelle ($s^{\mathrm{point}}_t$
élevé, $s^{\mathrm{context}}_t$ et $s^{\mathrm{collect}}_{t:t+K}$ faibles),
contextuelle ($s^{\mathrm{context}}_t$ élevé), collective
($s^{\mathrm{collect}}_{t:t+K}$ élevé), systémique (anomalie détectée dans
$m \geq 3$ pays simultanément).

\placeholder{Fig. G  Carte de confusion LOF/BiVAT par pays et type}{
Tableau 6×4. Lignes : Cameroun, Congo, Gabon, GQ, RCA, Tchad.
Colonnes : Ponctuelle, Contextuelle, Collective, Systémique.
Colormap : blanc (0\%) → vert foncé (100\% concordance).
Encadrés : cellules > 70\% concordance. Valeurs numériques dans chaque
cellule (\% concordance + nombre d'alertes entre parenthèses).
Fichier : \texttt{figG\_confusion\_lof\_bivat.pdf}
}

\subsubsection{Pertinence analytique}

Ce graphique synthétise la comparaison complète LOF--BiVAT en une seule
visualisation structurée selon deux dimensions analytiques indépendantes :
le pays (dimension géographique) et le type d'anomalie (dimension
typologique). Il répond à la question centrale de l'étude comparative :
sur quels types d'anomalies et dans quels pays les deux modèles s'accordent-ils,
et où divergent-ils ?

La structure attendue est révélatrice : forte concordance sur les anomalies
ponctuelles (les deux modèles détectent les mêmes outliers bruts) et faible
concordance sur les anomalies contextuelles et collectives (BiVAT en
détecte davantage grâce à sa modélisation temporelle). Si cette structure
apparaît clairement dans le graphique, elle valide empiriquement sur les
données BEAC l'argument théorique central des documents méthodologiques :
BiVAT n'est pas un LOF amélioré, c'est un outil qui détecte des types
d'anomalies structurellement différents.

\subsubsection{Positionnement dans le workflow}

\textbf{Phase~4  Synthèse comparative finale.} Ce graphique est le
dernier produit dans le workflow. Il est inclus dans la section de conclusion
du rapport opérationnel comme bilan de la comparaison inter-modèles et
guide de recommandation pour le déploiement.

\subsubsection{Données requises}

Classification de toutes les alertes LOF et BiVAT par type (ponctuelle,
contextuelle, collective, systémique) sur la fenêtre de test. Calcul du
taux de concordance par cellule (pays × type). Scores BiVAT multi-niveaux
$s^{\mathrm{point}}_t$, $s^{\mathrm{context}}_t$,
$s^{\mathrm{collect}}_{t:t+K}$ pour la classification.

\subsubsection{Décision supportée}

Bilan de la comparaison inter-modèles. Recommandation de déploiement :
LOF comme filtre rapide sur les anomalies ponctuelles, BiVAT comme outil
principal pour les anomalies contextuelles et collectives.
Documentation de la valeur ajoutée empirique de BiVAT sur les données
BEAC réelles. Guide pour l'allocation des ressources d'investigation.

\textbf{Priorité : Important.}

% ════════════════════════════════════════════════════════════════
\section{Catalogue Récapitulatif}

\begin{landscape}
\begin{longtable}{p{1.5cm}p{3.8cm}p{2cm}p{1.5cm}p{3cm}p{3.5cm}p{2cm}}
\caption{Catalogue complet des graphiques du système de visualisation.}
\label{tab:catalogue}\\
\toprule
\textbf{Fig.} & \textbf{Titre} & \textbf{Phase} &
\textbf{Modèle} & \textbf{Données requises} &
\textbf{Décision supportée} & \textbf{Priorité} \\
\midrule
\endfirsthead
\toprule
\textbf{Fig.} & \textbf{Titre} & \textbf{Phase} &
\textbf{Modèle} & \textbf{Données requises} &
\textbf{Décision supportée} & \textbf{Priorité} \\
\midrule
\endhead
\bottomrule\endlastfoot

Fig.~1 & Carte des données manquantes &
Ph.~1 & Commun &
Données brutes post-chargement &
Stratégie imputation ; exclusions &
Essentiel \\
\midrule

Fig.~2 & Heatmap corrélations résidus STL &
Ph.~1 & Commun &
Résidus RobustSTL + MAD, entraîn. &
Justification RPCA ; $k$ indicatif &
Essentiel \\
\midrule

Fig.~3 & Décomposition RobustSTL 4 panneaux &
Ph.~1 & Commun &
Brut + $T$, $S$, $R$ après RobustSTL &
Validation décomposition ; résidus &
Essentiel \\
\midrule

Fig.~A & Variance RPCA dans le temps &
Ph.~1 & Commun &
Résidus MAD + RPCA fenêtres 36m &
Ruptures régime structurelles ; $k$ &
Important \\
\midrule

Fig.~4 & Scree plot RPCA &
Ph.~2 & LOF &
Valeurs propres RPCA, entraîn. &
Choix de $k$ &
Essentiel LOF \\
\midrule

Fig.~5 & Biplot ACP (CP1 et CP2) &
Ph.~2 & LOF &
Projections CP1/CP2 + loadings &
Exploration visuelle préliminaire &
Complémentaire \\
\midrule

Fig.~6 & Projection UMAP &
Ph.~2 & LOF &
Résidus MAD + scores LOF* &
Exploration non-linéaire (annexe) &
Complémentaire \\
\midrule

Fig.~7 & Profils LOF(MinPts) &
Ph.~2 & LOF &
LOF$_k$(t) pour $k \in [10,30]$ &
Validation plage MinPts ; robustesse &
Essentiel LOF \\
\midrule

Fig.~8 & Distribution empirique LOF* &
Ph.~2 & LOF &
Scores LOF*, entraîn. &
Seuils $\tau_{\mathrm{IQR}}$ et $\tau_{95}$ &
Essentiel LOF \\
\midrule

Fig.~E & Intervalles CP BiVAT (calibration) &
Ph.~2 & BiVAT &
Scores BiVAT + CP, calib. set &
Validation coverage CP ; $\hat{q}_{1-\alpha}$ &
Essentiel BiVAT \\
\midrule

Fig.~9 & Boxplot LOF* par année &
Ph.~3 & LOF &
Scores LOF*, 2001--2026 &
Vue temporelle agrégée ; validation &
Important \\
\midrule

Fig.~10 & Série temporelle LOF*(t) &
Ph.~3 & LOF &
Scores LOF* + $\tau$ + événements &
Rapport opérationnel principal LOF &
Essentiel \\
\midrule

Fig.~11 & LOF + résidu + série brute &
Ph.~3 & LOF &
LOF*, $R_{i^*,t}$, $x_{i^*,t}$ &
Interprétation anomalie par anomalie &
Important \\
\midrule

Fig.~12 & Bar chart contributeurs &
Ph.~3 & LOF &
Résidus normalisés $R_{i,t}^*$ au mois $t$ &
Attribution indicateurs ; audit COBAC &
Essentiel \\
\midrule

Fig.~13 & Heatmap anomalies (ind. × mois) &
Ph.~3 & LOF &
$R_{i,t}^*$ pour tous mois anomaux &
Interprétation économique globale &
Essentiel \\
\midrule

Fig.~14 & Carte de chaleur inter-pays &
Ph.~3 & LOF &
LOF* 6 pays + seuils nationaux &
Anomalies systémiques CEMAC &
Essentiel \\
\midrule

Fig.~B & Concordance Actif--Passif &
Ph.~3 & LOF &
LOF* Actif + LOF* Passif par pays &
Triage événement réel vs collecte &
Essentiel \\
\midrule

Fig.~D & Comparaison LOF vs BiVAT &
Ph.~4 & Commun &
LOF* + BiVAT normalisés + seuils &
Alertes niveau~1 vs niveau~2 &
Essentiel \\
\midrule

Fig.~E bis & Score BiVAT + intervalles CP &
Ph.~4 & BiVAT &
Scores BiVAT + CP, test &
Rapport opérationnel principal BiVAT &
Essentiel \\
\midrule

Fig.~F & Attribution SHAP top-5 BiVAT &
Ph.~4 & BiVAT &
Valeurs DeepSHAP $\phi_d$ aux top-5 &
Attribution réglementaire COBAC &
Essentiel BiVAT \\
\midrule

Fig.~G & Carte confusion LOF/BiVAT &
Ph.~4 & Commun &
Alertes classifiées par type + pays &
Bilan comparatif ; recommandation &
Important \\
\end{longtable}
\end{landscape}

% ════════════════════════════════════════════════════════════════
\section{Conclusion et Guide d'Utilisation Opérationnelle}

\subsection{Bilan du système visuel}

Ce document a décrit un système de visualisation complet en 21 graphiques
organisés en quatre phases séquentielles, couvrant l'intégralité du workflow
de détection de points d'intérêt dans les données monétaires BEAC-CEMAC.
Parmi les 14 graphiques du corpus initial, 12 sont retenus, 2 sont reclassés
en complémentaires avec des réserves méthodologiques documentées (Fig.~5
biplot ACP et Fig.~6 projection UMAP), et 1 est corrigé (Fig.~12 : radar
chart remplacé par bar chart horizontal pour des raisons d'auditabilité).
Sept graphiques supplémentaires ont été identifiés comme nécessaires à la
complétude du système (Fig.~A, Fig.~B, Fig.~D, Fig.~E, Fig.~E bis,
Fig.~F, Fig.~G).

\subsection{Guide de séquencement opérationnel}

Le tableau~\ref{tab:sequencement} présente le guide d'utilisation
opérationnelle du système visuel, organisé par ordre d'exécution et par
question analytique à laquelle chaque graphique répond.

\begin{table}[ht]
\centering
\caption{Guide de séquencement opérationnel des graphiques.}
\label{tab:sequencement}
\small
\begin{tabular}{p{0.5cm}p{1.8cm}p{2cm}p{4cm}p{5cm}}
\toprule
\textbf{\#} & \textbf{Figure} & \textbf{Quand} &
\textbf{Question posée} & \textbf{Si anomalie détectée} \\
\midrule
1 & Fig.~1 & Phase~1, t=0 &
Les données sont-elles exploitables ? &
Appliquer la stratégie d'imputation adaptée \\
\midrule
2 & Fig.~3 & Phase~1, après STL &
La décomposition est-elle correcte ? &
Ajuster paramètres RobustSTL si résidus non stationnaires \\
\midrule
3 & Fig.~2 & Phase~1, après STL &
Quelle est la structure de covariance des résidus ? &
Informer le choix de $k$ pour la RPCA \\
\midrule
4 & Fig.~A & Phase~1, après STL &
La structure de covariance change-t-elle dans le temps ? &
Activer CP pondérée pour les périodes à faible variance expliquée \\
\midrule
5 & Fig.~4 & Phase~2, LOF &
Combien de composantes RPCA retenir ? &
Fixer $k$ au coude ou à 90\,\% de variance expliquée \\
\midrule
6 & Fig.~7 & Phase~2, LOF &
La plage MinPts $[10,30]$ est-elle valide ? &
Ajuster les bornes si aucune anomalie n'est robuste \\
\midrule
7 & Fig.~8 & Phase~2, LOF &
Quels sont les seuils de détection $\tau$ ? &
Fixer $\tau_{\mathrm{IQR}}$ et $\tau_{95}$ \\
\midrule
8 & Fig.~E & Phase~2, BiVAT &
Le coverage CP est-il conforme au niveau nominal ? &
Ajuster SSBC si écart $>5\,\%$ \\
\midrule
9 & Fig.~10 & Phase~3, LOF &
Quand les anomalies LOF se produisent-elles ? &
Annoter avec événements CEMAC ; investiguer pics non annotés \\
\midrule
10 & Fig.~E bis & Phase~3, BiVAT &
Quand les anomalies BiVAT se produisent-elles ? &
Classer par certitude : score + intervalle inférieur vs seuil \\
\midrule
11 & Fig.~13 & Phase~3, LOF &
Quels indicateurs portent chaque anomalie LOF ? &
Produire Fig.~12 (bar chart) pour les anomalies robustes \\
\midrule
12 & Fig.~F & Phase~3, BiVAT &
Quels indicateurs portent chaque anomalie BiVAT ? &
Documentation réglementaire COBAC \\
\midrule
13 & Fig.~14 & Phase~3, commun &
Les anomalies sont-elles systémiques CEMAC ? &
Si $m \geq 3$ pays : alerte systémique niveau~1 \\
\midrule
14 & Fig.~B & Phase~3, LOF &
L'anomalie est-elle concordante Actif--Passif ? &
Concordante : événement économique réel. Discordante : erreur collecte \\
\midrule
15 & Fig.~D & Phase~4, commun &
LOF et BiVAT s'accordent-ils sur cette anomalie ? &
Accord : alerte niveau~1. LOF seul : anomalie ponctuelle.
BiVAT seul : anomalie contextuelle/collective \\
\midrule
16 & Fig.~G & Phase~4, commun &
Quel modèle est le plus pertinent pour quel type d'anomalie ? &
Recommandation de déploiement : LOF filtrage rapide,
BiVAT surveillance contextuelle \\
\bottomrule
\end{tabular}
\end{table}

\subsection{Recommandation de déploiement}

Sur la base du système visuel documenté, la recommandation opérationnelle
pour la BEAC est la suivante.

\textit{Usage mensuel en surveillance continue.} Les graphiques Fig.~10
(série temporelle LOF), Fig.~E bis (score BiVAT + intervalles CP) et
Fig.~14 (heatmap inter-pays) sont mis à jour mensuellement à chaque
réception des nouvelles données de reporting bancaire. Ils constituent
le tableau de bord de surveillance de premier niveau.

\textit{Investigation de second niveau.} Pour chaque nouvelle alerte
identifiée (score LOF* ou BiVAT dépassant leurs seuils respectifs), les
graphiques Fig.~11 (superposition LOF + résidu + brut), Fig.~12 (bar chart
contributeurs LOF), Fig.~F (attribution SHAP BiVAT) et Fig.~B (concordance
Actif--Passif) sont produits pour cette alerte spécifique. Ils constituent
le dossier d'investigation transmis à l'équipe de surveillance.

\textit{Rapport annuel de bilan.} Les graphiques Fig.~9 (boxplot par année),
Fig.~13 (heatmap globale des anomalies) et Fig.~G (carte de confusion
LOF/BiVAT) sont produits annuellement comme bilan de la surveillance et
comme outil d'amélioration continue des deux modèles.

\textit{Graphiques internes de validation.} Les graphiques Fig.~4 (scree
plot), Fig.~5 (biplot), Fig.~6 (UMAP), Fig.~7 (profils MinPts) et Fig.~8
(distribution LOF*) sont des outils de calibration interne, produits lors
de la mise en place initiale et lors de chaque recalibration annuelle des
modèles. Ils ne sont pas inclus dans les rapports opérationnels destinés
aux superviseurs BEAC.

% ════════════════════════════════════════════════════════════════
\section*{Bibliographie}
\addcontentsline{toc}{section}{Bibliographie}

\begin{enumerate}

\item Breunig, M.M., Kriegel, H.-P., Ng, R.T., \& Sander, J. (2000).
LOF: Identifying Density-Based Local Outliers. \textit{ACM SIGMOD 2000},
pp.~93--104. DOI:~10.1145/342009.335388.

\item Wen, Q. et al. (2019). RobustSTL: A Robust Seasonal-Trend
Decomposition Algorithm for Long Time Series. \textit{AAAI-19}.
arXiv:1812.01767.

\item Candès, E.J., Li, X., Ma, Y., \& Wright, J. (2011). Robust Principal
Component Analysis? \textit{Journal of the ACM}, 58(3), Article~11.
DOI:~10.1145/1970392.1970395.

\item Zimek, A., Schubert, E., \& Kriegel, H.-P. (2012). A survey on
unsupervised outlier detection in high-dimensional numerical data.
\textit{Statistical Analysis and Data Mining}, 5(5), 363--387.

\item Chandola, V., Banerjee, A., \& Kumar, V. (2009). Anomaly Detection:
A Survey. \textit{ACM Computing Surveys}, 41(3), Article~15.

\item Xu, J., Wu, H., Wang, J., \& Long, M. (2022). Anomaly Transformer:
Time Series Anomaly Detection with Association Discrepancy.
\textit{ICLR 2022 (Spotlight)}.

\item Lundberg, S.M., \& Lee, S.I. (2017). A Unified Approach to
Interpreting Model Predictions. \textit{NeurIPS}, 30.

\item Kim, T.S. et al. (2022). Reversible Instance Normalization for
Accurate Time-Series Forecasting against Distribution Shift.
\textit{ICLR 2022}.
\url{https://openreview.net/forum?id=cGDAkQo1C0p}

\item Vovk, V., Gammerman, A., \& Shafer, G. (2005).
\textit{Algorithmic Learning in a Random World}. Springer.

\item arXiv:2509.15349 (2025). Small Sample Beta Correction for Conformal
Prediction.

\item Xu, C., \& Xie, Y. (2021). Conformal Prediction Intervals for Dynamic
Time-Series. arXiv:2010.09107.

\item Tukey, J.W. (1977). \textit{Exploratory Data Analysis}.
Addison-Wesley.

\item BIS / FSI (2024). How regulators can address AI explainability.
\url{https://www.bis.org/fsi/fsipapers24.pdf}

\item ECB (2025). Guide to Internal Models (ML chapter).
\url{https://www.bankingsupervision.europa.eu}

\item SHARC (2025). SHAP-Based Interpretability in ML Risk Models for
ICAAP and CCAR. arXiv:2607.05484.

\item arXiv:2403.04429 (2024). Exploring the Influence of Dimensionality
Reduction on Anomaly Detection in MTS.

\item arXiv:2508.21128 (2026). STAD: Seasonal-Trend Anomaly Detection.

\item Goldstein, M., \& Uchida, S. (2016). A Comparative Evaluation of
Unsupervised Anomaly Detection Algorithms for Multivariate Data.
\textit{PLoS ONE}, 11(4): e0152173.

\item Moukoko Ekambi, G.M. (2025). Méthodologie Modèle~1 : Détection
d'Anomalies par LOF sur les Données Monétaires BEAC. Rapport de stage,
BEAC / ENSPD / Université de Douala, LIDIA Laboratory.

\item Moukoko Ekambi, G.M. (2025). Détection de Variations Atypiques dans
les Séries Temporelles Financières par Architectures Hybrides VAE-Transformer :
État de l'art outillé et justification architecturale de BiVAT. Rapport de
stage, BEAC / ENSPD / Université de Douala, LIDIA Laboratory.

\item Moukoko Ekambi, G.M. (2025). Pipeline de prétraitement des données
monétaires BEAC-CEMAC : étapes communes et étapes spécifiques aux Modèles
LOF et BiVAT. Rapport de stage, BEAC / ENSPD / Université de Douala,
LIDIA Laboratory.

\end{enumerate}

\vspace{2cm}
\hrule
\vspace{0.5cm}
\noindent{\small\textit{Document rédigé dans le cadre du stage de recherche
à la Banque des États de l'Afrique Centrale (BEAC), M2 Data Science \& IA,
ENSPD / Université de Douala  LIDIA Laboratory. Août 2026.}}

\end{document}